\documentclass[10pt]{scrartcl}
\usepackage[mdthm,colorsec,noauthor,noasy]{evan}
\usepackage{handout}
\DeclareMathOperator{\im}{range}\DeclareMathOperator{\tr}{tr}
\title{MATH 110 --- Week 10 --- Solutions}
\subtitle{Notes Chapter 18; Homework 10}
\author{}\date{}\ihead{}
\begin{document}\maketitle
\begin{prob}\textbf{HW10 P1.}\par\idea{Subtract earlier polynomial components; parity makes the odd mixed integrals zero.}\begin{steps}\ii Start with $u_1=1$, $u_2=x$ and $u_3=x^2-\langle x^2,1\rangle/\|1\|^2=x^2-1/3$. Parity gives $\langle x,1\rangle=\langle x^2,x\rangle=0$. Also $\langle x^2,1\rangle=\int_{-1}^1x^2\,dx=2/3$ and $\|1\|^2=2$, giving $1/3$.
\ii The squared norms are $\|1\|^2=2$, $\|x\|^2=2/3$, and
\[\|x^2-1/3\|^2=\int_{-1}^1(x^4-2x^2/3+1/9)\,dx=2/5-4/9+2/9=8/45.\]\ii Divide by those norms' positive square roots to obtain $e_1=1/\sqrt2$, $e_2=\sqrt{3/2}\,x$, $e_3=\sqrt{5/8}(3x^2-1)$. Proposition 87 gives an orthonormal basis with the original initial spans.\end{steps}\end{prob}
\begin{prob}\textbf{HW10 P2.}\par\idea{Orthogonality to a spanning list is a pair of homogeneous equations.}\begin{steps}\ii For $(a,b,c,d)$ the equations are $a-b+c=0$ and $2b+c=0$. Thus $c=-2b$ and $a=3b$, while $b,d$ are free.\ii Therefore $W^\perp=\operatorname{span}((3,1,-2,0),(0,0,0,1))$. The displayed vectors are independent by their second and fourth coordinates, so form a basis.\end{steps}\end{prob}
\begin{prob}\textbf{HW10 P3.}\par\idea{Use one integral equation and choose two independent solutions.}\begin{steps}\ii For $p=a+bx+cx^2$, $\langle p,x+1\rangle=2a+2b/3+2c/3$. Its vanishing means $a=-(b+c)/3$.\ii Thus $p=b(x-1/3)+c(x^2-1/3)$. These two polynomials are independent by their $x$ and $x^2$ coefficients, so $x-1/3,x^2-1/3$ is a basis of $W^\perp$.\end{steps}\end{prob}
\begin{prob}\textbf{HW10 P4.}\par\idea{Take one orthogonal complement of a sum, then take the complement again.}\begin{steps}\ii The sum identity gives $(U^\perp+W^\perp)^\perp=(U^\perp)^\perp\cap(W^\perp)^\perp=U\cap W$.\ii Take complements of both sides. Every subspace is finite-dimensional here, so the double-complement theorem gives $U^\perp+W^\perp=(U\cap W)^\perp$.\end{steps}\end{prob}
\begin{prob}\textbf{HW10 P5.}\par\idea{The orthogonal residual contributes the same squared error for every candidate in the subspace.}\begin{steps}\ii Decompose $v=u_0+w$ with $u_0\in U,w\in U^\perp$. For each $j$, $\langle v,e_j\rangle=\langle u_0,e_j\rangle$, so Proposition 86 gives $P_Uv=u_0=\sum_j\langle v,e_j\rangle e_j$.\ii For any $u\in U$, $v-u=(v-P_Uv)+(P_Uv-u)$ is an orthogonal sum. Hence $\|v-u\|^2=\|v-P_Uv\|^2+\|P_Uv-u\|^2$.\ii This proves the inequality, with equality exactly when $P_Uv-u=0$, that is $u=P_Uv$.\end{steps}\end{prob}
\begin{prob}\textbf{HW10 P6 (optional).}\par\idea{A functional on $\mathcal P_1$ is determined by its values on $1,x$.}\begin{steps}\ii Write $p=a+bx$. Testing $f=1$ gives $2=\int_0^21\,dx=\langle1,p\rangle=2a$, so $a=1$.\ii Testing $f=x$ gives $2=\int_0^2x\,dx=\langle x,p\rangle=2b/3$, so $b=3$. Thus $p=1+3x$.\ii Linearity verifies the identity for every $f=c+dx$. Riesz uniqueness, or uniqueness of the two coefficient equations, proves uniqueness.\end{steps}\end{prob}
\begin{prob}\textbf{HW10 P7 (optional).}\par\idea{The norm bound forces the projector's range and kernel to be orthogonal, by testing one mixed scalar.}\begin{steps}\ii Orthogonal projection has the norm bound by Pythagoras, as proved in the projection-properties theorem.\ii Conversely assume the bound. Proposition 91 gives $V=\im P\oplus\ker P$. For $u\in\im P,w\in\ker P$, $P(u+tw)=u$ for every scalar $t$.\ii Put $c=\langle u,w\rangle$. If $w=0$ then $c=0$. Otherwise choose $t=-c/\|w\|^2$. Expanding gives $\|u+tw\|^2=\|u\|^2-|c|^2/\|w\|^2$.\ii The assumed norm bound requires this to be at least $\|u\|^2$, so $c=0$. Hence $\ker P\subseteq(\im P)^\perp$. Their dimensions are both $\dim V-\dim\im P$, by rank-nullity and Proposition 89, so they are equal.\ii Therefore $P$ projects onto its range along its orthogonal complement, which is precisely orthogonal projection. \end{steps}\end{prob}
\begin{prob}\textbf{LADR 6A Exercise 11.}\par\idea{Find the line component by its single inner product.}\begin{steps}\ii For $w=(1,3)$ and $a=(1,2)$, $\langle a,w\rangle=7$ and $\|w\|^2=10$. Thus $u=7w/10=(7/10,21/10)$.\ii Subtract to get $v=a-u=(3/10,-1/10)$, whose dot product with $w$ is $3/10-3/10=0$. This gives all requested properties; uniqueness follows from the orthogonal direct sum.\end{steps}\end{prob}
\begin{prob}\textbf{LADR 6A Exercise 6.}\par\idea{A mixed inner product can be eliminated by a scalar choice, forcing a strict decrease unless it is zero.}\begin{steps}\ii If $\langle u,v\rangle=0$, expanding gives $\|u+av\|^2=\|u\|^2+|a|^2\|v\|^2\ge\|u\|^2$.\ii Conversely $v=0$ makes the conclusion immediate. For $v\ne0$, set $a=-\langle u,v\rangle/\|v\|^2$. Then $\|u+av\|^2=\|u\|^2-|\langle u,v\rangle|^2/\|v\|^2$.\ii The bound for this $a$ forces the subtracted term to be zero, so $\langle u,v\rangle=0$.\end{steps}\end{prob}
\end{document}
